% ===================================================================
%  TABELL Nr. 10 — D'Mier. Den Hafen.
%  Traduction 2026 d'apres texto/io, controlee sur texto/fr, texto/de
%  et texto/nl. Consigne : voir texto/lb/10-tabell-01.tex.
%
%  LE LUXEMBOURG N'A PAS DE MER, ET C'EST LA SITUATION TCHEQUE. La
%  colonne treize avait montre ce qu'une langue fait alors : elle
%  prete a la mer les mots de la montagne — « útes » pour la falaise
%  et pour l'ecueil. LE LUXEMBOURGEOIS NE FAIT PAS CELA. Il prend le
%  vocabulaire tout entier a l'allemand, et sans une hesitation :
%  « Schëff », « Hafen », « Anker », « Segel », « Mast »,
%  « Matrous », « Klippe », « Rëff », « Brandung », « Gezäiten ».
%
%  ET C'EST LA DIFFERENCE ENTRE MANQUER D'UNE CHOSE ET MANQUER D'UNE
%  CHOSE QU'UN VOISIN A. Le tcheque n'avait pas de mer et n'avait pas
%  non plus de voisin dont la langue lui fut assez proche pour qu'il
%  puisse s'y servir ; il a donc fabrique. Le luxembourgeois n'a pas
%  de mer, mais il a un voisin qui en a une ET dont la langue est la
%  sienne a un cheveu pres : il n'a rien eu a fabriquer. UN MANQUE NE
%  PRODUIT PAS TOUJOURS UNE INVENTION ; IL FAUT ENCORE QU'IL N'Y AIT
%  PERSONNE A COTE.
%
%  MAIS LA RIVIERE, ELLE, EST A LUI. Le pays a la Moselle et la Sure,
%  et le batelier de l'alinea 15 s'y dit en luxembourgeois :
%  « Floss », « Baach », « Ufer », « Boot », « Rudder »,
%  « Fëscher », « Schlepper », « Kai ». La ligne du tableau 8 —
%  riche ou le pays est riche — passe ici entre la riviere et la mer,
%  a l'interieur du meme tableau.
%
%  ET LES GRADES DE MARINE SONT EMPRUNTES EN BLOC, ce qui n'etonnera
%  pas : « Leitnant », « Kommandant », « Admiral »,
%  « Quartiermeeschter ». Le pays a une armee et pas de flotte.
% ===================================================================

%%K t10-apar-1 apar
\VUsaut{4.0mm}
\VUtitre{15.0pt}{60}{TABELL Nr. 10}
\VUsaut{3.0mm}
\VUpk{11.4pt}{40}{D'Mier. --- Den Hafen.}
\VUsaut{3.0mm}

%%K t10-01-1 p
1. --- Am Summer, während déi eng Rou a Killt um Bierg sichen, ginn déi aner léiwer
un d'Küst. Esou hu mir et d'lescht Joer gemaach, well mir hunn den \VUgras{Ozean}
\textsuperscript{(1)} ganz gär.

%%K t10-02-1 p
2. --- Wat mech ubelaangt, hunn ech e grousst Vergnügen, déi immens Waasserweit ze
betruechten, déi sech bis un den \VUgras{Horizont} \textsuperscript{(2)} zitt, an déi
enorm \VUgras{Dampschëffer} \textsuperscript{(3)} fortfuere ze gesinn, op déi
d'Reesend ariichten, déi op Amerika ginn.

%%K t10-03-1 p
3. --- Dofir wielt mäi Papp, dee meng Virléift kennt, fir d'Vakanz ëmmer eng ganz
roueg Plage, déi nieft engem vun eise groussen \VUgras{Krichshäfen} läit.

%%K t10-03-2 p
Bei eisem leschten Openthalt ass d'\VUgras{Escadre} hannert dem enorme \VUgras{Damm}
\textsuperscript{(4)} viru Anker gaangen, deen de \VUgras{Krichshafen}
\textsuperscript{(5)} zoumécht a schützt, an et ass mer gelongen, no mengem Wonsch
déi verschidde \VUgras{Krichsschëffer} ze bewonneren, dat klengt
\VUgras{Ënnerseeboot} \textsuperscript{(10)}, dat ënner d'Waasser tauche a
verschwannt, an och dat enormt \VUgras{Panzerschëff} \textsuperscript{(9)}, dat bis
elo bal nëmmen d'Attacke vum \VUgras{Torpedoboot} \textsuperscript{(8)} gefaart huet.

%%K t10-tit-2 sub
\VUsaut{3.0mm}
\VUpk{10.2pt}{40}{Déi kleng Schëffer.}
\VUsaut{3.0mm}

%%K t10-04-1 p
4. --- Dëst wousst schonn ganz \VUgras{geschéckt} de schwaache Punkt vum décke
Metallpanzer ze fannen, fir do den geféierlechen \VUgras{Torpedo} anzestiechen, deem
seng Explosioun de Verloscht vum Schëff verursaacht; mä trotzdeem war et manner ze
faarten wéi d'Ënnerseeboot, well d'Zerstéierer verfollegen et liicht.

%%K t10-05-1 p
5. --- Ech hunn och déi séier gepanzert \VUgras{Croiseuren} \textsuperscript{(6)}
gesinn, bal esou onverwundbar wéi e richtegt Panzerschëff, e puer
\VUgras{Transportschëffer} \textsuperscript{(12)}, dräi oder véier
\VUgras{Kanounebooter} \textsuperscript{(7)} a schlussendlech eng \VUgras{Fregatt}
\textsuperscript{(11)}. \VUgras{D'Schauspill vun där Flott, wa si manövréiert fir
d'Anker ze liichten, ass héchst interessant.}

%%K t10-06-1 p
6. --- Wéi e Bau-Ënnerscheed tëscht deene grousse Stolmassen an den elegante
\VUgras{Dräimasteren} oder de zierleche \VUgras{Segelschëffer}
\textsuperscript{(13)}, déi nach haut an der Handelsflott benotzt ginn, déi ech an den
Hafen erafuere gesinn hunn, mat alle Segelen, wéi ech um \VUgras{Kai}
\textsuperscript{(14)} spadséiert sinn. Wierklech hunn dës vill vun hire Virdeeler
verluer, \VUgras{wéi} de \VUgras{Wand} géint si war. Si hu missen all d'Segelen
erofloossen, fir nees an de Baseng ze kommen, an e \VUgras{Schlepper}
\textsuperscript{(16)} war néideg, fir se ze sichen a se, wéi einfach
\VUgras{Lastbooter} \textsuperscript{(17)}, un de Kai ze bréngen.

%%K t10-tit-3 sub
\VUsaut{3.0mm}
\VUpk{10.2pt}{40}{Den Handelshafen.}
\VUsaut{3.0mm}

%%K t10-07-1 p
7. --- Wat un deem Hafen interessant ass, vun deem ech schwätzen, ass, datt en net
nëmmen e Krichshafen huet, dee kënschtlech mat risege Aarbechte geriicht gouf, mä och
e wonnerbaren \VUgras{Handelshafen,} deen an der \VUgras{Mündung} vun engem grousse
Floss läit.

%%K t10-08-1 p
8. --- D'Landzong, déi déi zwee Basenge trennt, ass befestegt a gëtt vun enger Rei
\VUgras{Batterien} an \VUgras{Bastiounen} verdeedegt, déi an d'Mier virstoussen. Ee
brauch eng speziell Erlaabnes, fir op de \VUgras{Wäll} \textsuperscript{(15)} ze
spadséieren. Ech hu se liicht duerch mäi Monni kritt, deen Ingenieur vun de
\VUgras{Werften} \textsuperscript{(18)} ass, an ech sinn d'Schëffer kucke gaangen,
déi ee ënner grousse Hallen baut.

%%K t10-09-1 p
9. --- Ech war esouguer beim Stapellaf vun engem Panzerschëff dobäi, an ech erënnere
mech un dee bewegenden Ablack, deen déi ganz Masse gespuert huet, wéi déi risech
Mass, sech selwer iwwerlooss, op hirem Wéirahmen ugefaangen huet ze glidden an um
Waasser vum Floss schwammen ass.

%%K t10-10-1 p
10. --- Fir bei d'Werften ze kommen, geet ee duerch dat \VUgras{Manöverfeld}
\textsuperscript{(19)}, wou d'Zaldoten aus der Garnisoun all Dag üben kommen, an ee
geet iwwer en eisenen \VUgras{Steg} \textsuperscript{(20)}, deen d'Paradenplaz mam
\VUgras{Arsenal} \textsuperscript{(21)} verbënnt.

%%K t10-tit-4 sub
\VUsaut{3.0mm}
\VUpk{10.2pt}{40}{D'Arsenal an d'Docken.}
\VUsaut{3.0mm}

%%K t10-11-1 p
11. --- Mat mengem Monni hunn ech dee grousse Bau besicht, voller \VUgras{Kanounen}
\textsuperscript{(22)}, \VUgras{Kugelen} \textsuperscript{(23)} an \VUgras{Obussen.}
Ech hunn och d'\VUgras{Metallgisserei} \textsuperscript{(24)} gesinn, an där
d'\VUgras{Kessel} \textsuperscript{(25)} vun den Dampschëffer gemaach ginn.

%%K t10-12-1 p
12. --- Ganz nobäi sinn d'\VUgras{Docken} \textsuperscript{(26)} --- d'Reparaturdocken
--- an déi ganz bemierkenswäert \VUgras{schwiewend Docken} \textsuperscript{(27)}, an
deenen ech ganz no bei engem Schëff, dat ze reparéiere war, d'\VUgras{Schrauf}
\textsuperscript{(28)}, de \VUgras{Kiel} \textsuperscript{(29)} an dat
\VUgras{Rudder} \textsuperscript{(30)} vun engem Boot gesi konnt.

%%K t10-13-1 p
13. --- Den Handelshafen ass genau esou studéierenswäert wéi de Krichshafen, an ech
hu vill Stonne gapend op de Kaie verbruecht, wou d'\VUgras{Raddampfer}
\textsuperscript{(31)} ugebonne sinn, déi de Floss erop fueren, a wou ech
d'\VUgras{Mastkrunn} \textsuperscript{(32)} funktionéiere gesinn hunn, déi wéi
Fiedere risech Laaschten hiewen.

%%K t10-tit-5 sub
\VUsaut{4.0mm}
\VUpk{10.2pt}{40}{De Lotse.}
\VUsaut{3.0mm}

%%K t10-14-1 p
14. --- Ech hunn d'\VUgras{Transatlantiker} \textsuperscript{(34)} bewonnert, déi aus
der Rieds erausfueren, mat hire \VUgras{Fändelen} \textsuperscript{(35)} an der
Loft, gefouert vun engem geschéckte \VUgras{Lotse} \textsuperscript{(36)}, dee
wonnerbar der \VUgras{Fahrrënn} ze follege weess, déi mat \VUgras{Bojen}
\textsuperscript{(40)} an \VUgras{Baken} markéiert ass, an deen de
\VUgras{Lastkähn} \textsuperscript{(41)} ausweicht, déi mat hirem eenzege
véiereckege Segel schwéier ze steiere sinn. Ech hunn um \VUgras{Virschëff}
\textsuperscript{(37)} --- dem Bug --- d'\VUgras{Kabinnen} \textsuperscript{(39)} vun
der zweeter Klass erkannt, an um \VUgras{Hannerschëff} \textsuperscript{(38)} ---
dem Heck --- d'Salongen fir d'Passagéier vun der éischter Klass.

%%K t10-15-1 p
15. --- Heiansdo war ech och beim Lueden vun engem \VUgras{Frachtschëff}
\textsuperscript{(42)} dobäi, an dat grad d'Wuere vum \VUgras{Lager}
\textsuperscript{(43)} a vun der \VUgras{Séistatioun} \textsuperscript{(44)}
opgehäift goufen, déi mat de villen Zich vun der \VUgras{Kaibunn}
\textsuperscript{(45)} bruecht goufen.

%%K t10-tit-6 sub
\VUsaut{3.0mm}
\VUpk{10.2pt}{40}{D'Rudervergnügen.}
\VUsaut{3.0mm}

%%K t10-16-1 p
16. --- Dacks hunn ech am \VUgras{Schwämmbaseng} \textsuperscript{(53)} gerudert, an
deen d'\VUgras{Barken} \textsuperscript{(46)} sech flüchten, an et war eng Freed fir
mech, dat \VUgras{Rudder} \textsuperscript{(47)} ze féieren. Ech si bis op dat
\VUgras{Deck} \textsuperscript{(48)} vun enger \VUgras{Segelbark}
\textsuperscript{(49)} geklommen, ech sinn u \VUgras{Tauen} \textsuperscript{(52)}
um \VUgras{Mast} \textsuperscript{(51)} bis an d'\VUgras{Rahen}
\textsuperscript{(50)} geklommen, duerno hunn ech mäi Spadséiergang \VUgras{mat enger
Yol} \textsuperscript{(54)} tëscht de schwéiere \VUgras{Fëscherbooter}
\textsuperscript{(55)} weidergefouert, wou d'\VUgras{Netzer} \textsuperscript{(56)}
dréchnen.

%%K t10-17-1 p
17. --- Um \VUgras{Kai} \textsuperscript{(58)} sinn déi verschiddenste
\VUgras{Wueren} \textsuperscript{(59)} viru mir laanschtgaangen, a \VUgras{Këschten}
\textsuperscript{(60)} oder a \VUgras{Ballen} \textsuperscript{(61)}. Si hunn
d'\VUgras{Ladung} \textsuperscript{(63)} vun de Lastbooter gebilt, a mächteg
\VUgras{Krunnen} \textsuperscript{(62)} hu se erofgeholl an op \VUgras{Camionen}
\textsuperscript{(64)} gesat, während d'\VUgras{Fuerleit} se deklaréiert an
d'Taxen am \VUgras{Zollamt} \textsuperscript{(65)} bezuelt hunn.

%%K t10-18-1 p
18. --- Schlussendlech, wéi ech bei d'\VUgras{Dréibréck} \textsuperscript{(66)} oder
bei d'\VUgras{Schleis} \textsuperscript{(69)} komm sinn, virum \VUgras{Bagger}
\textsuperscript{(68)} laanschtgaangen, deen de \VUgras{Kanal} \textsuperscript{(67)}
botzt, sinn ech um Kai nieft dem \VUgras{Semaphor} \textsuperscript{(70)}
ausgeklommen.

%%K t10-19-1 p
19. --- Op der anerer Säit vun deem Kai leeën déi \VUgras{Schaluppen}
\textsuperscript{(71)} un, déi d'Offizéier vun der Escadre un d'Land bréngen. Et ass
e bewonnerenswäert Schauspill, ze gesinn, wéi d'Matrousen am Takt hir Rudder hiewen,
während d'Bark, vun der \VUgras{Well} \textsuperscript{(72)} gedroen, liicht bei der
Trap vum Ausklammplaz uleet. Op där Plaz kann ee besser d'\VUgras{Gezäiten} an
d'Bewegung vun der \VUgras{Flut} an der \VUgras{Ebb} op der \VUgras{Plage}
\textsuperscript{(73)} beobachten.

%%K t10-tit-7 sub
\VUsaut{3.0mm}
\VUpk{10.2pt}{40}{De Club.}
\VUsaut{3.0mm}

%%K t10-20-1 p
20. --- Fir heemzekommen, hunn ech den \VUgras{elektreschen Tram}
\textsuperscript{(74)} geholl, deem säin \VUgras{Arrêt} \textsuperscript{(75)} nieft
dem \VUgras{Poteau} ass, deen virum Offizéiersclub steet.

%%K t10-20-2 p
Vum \VUgras{Balkon} \textsuperscript{(76)} vum Club, dee mat \VUgras{Ankeren}
\textsuperscript{(77)}, Kanounen a Kugele geschmückt ass, hunn ech dacks eng
herrlech Versammlung vu Marineoffizéiere gesinn: \VUgras{Leitnanten}
\textsuperscript{(78)} mat hirem Degen, \VUgras{Artilleriekapitänen}
\textsuperscript{(79)} an \VUgras{Infanteriekapitänen} \textsuperscript{(80)} mat
hirem \VUgras{Säbel.} Hannert hinne stoung e \VUgras{Matrous} \textsuperscript{(81)},
deen hinnen e \VUgras{Fernrouer} bruecht huet, wa si wollte gesinn, wat wäit ewech
geschitt.

%%K t10-21-1 p
21. --- Ech hunn och jonk \VUgras{Ënnerleitnanten} \textsuperscript{(82)} gesinn, déi
d'Instruktioune vun hirem \VUgras{Kommandant} \textsuperscript{(83)} oder vum
\VUgras{Admiral} \textsuperscript{(84)} kritt hunn, während e
\VUgras{Quartiermeeschter} \textsuperscript{(85)} gewaart huet, fir se op d'Schëff ze
bréngen.

%%K t10-22-1 p
22. --- Vun deem Balkon ass d'Aussiicht herrlech: ee beherrscht déi ganz Plage mat
hiren \VUgras{Zelter} \textsuperscript{(86)} an \VUgras{Kabinnen}
\textsuperscript{(87)}, an ee gesäit, zur Stonn vum Bueden, d'\VUgras{Bueder}
\textsuperscript{(88)}, déi frou virum \VUgras{Casino} \textsuperscript{(89)}
dollen.

%%K t10-23-1 p
23. --- Um Enn vun der Plage bilt d'Küst eng kleng \VUgras{fielseg}
\VUgras{Hallefinsel} \textsuperscript{(91)}, wou d'\VUgras{Brandungswell}
\textsuperscript{(92)} sech brécht a wou e \VUgras{Luuchttuerm} \textsuperscript{(93)}
gebaut ass, deen nuets de Schëffer de Wee weist. Déi Hallefinsel ass mam Land nëmmen
duerch eng ganz schmuel \VUgras{Landeng} \textsuperscript{(90)} verbonnen.

%%K t10-tit-8 sub
\VUsaut{3.0mm}
\VUpk{10.2pt}{40}{Allgemengt Bild vun de Küsten.}
\VUsaut{3.0mm}

%%K t10-24-1 p
24. --- D'\VUgras{Klipp} \textsuperscript{(94)} zitt sech laanscht, vum Mier
gespléckt, dat launesch dran eng Rei \VUgras{Buchten} \textsuperscript{(95)} an
\VUgras{Kappen} zeechent, déi se héchst molerësch maachen.

%%K t10-24-2 p
Um \VUgras{Plateau} \textsuperscript{(96)}, dat den \VUgras{Mieresarm}
\textsuperscript{(98)} beherrscht, ass e ganz wichtegt \VUgras{Fort}
\textsuperscript{(97)} an e puer „Villaen“.

%%K t10-25-1 p
25. --- Deen Arm trennt vum Festland eng ganz geféierlech \VUgras{Insel}
\textsuperscript{(99)}, déi vu \VUgras{Rëffer} \textsuperscript{(100)} an
\VUgras{Brandungen} ëmginn ass, wou Barken a souguer grouss Schëffer heiansdo
stranden. Trotz der Séiergkeet vun der Hëllef an der schneller Arrivée vun de
\VUgras{Rettungsbooter} sinn déi \VUgras{Schëffsbréch} schrecklech, a während de
\VUgras{Äquinoktialstierm} sinn e puer Schëffer mat Mann a Wuer ënnergaangen.

%%K t10-25-2 p
Trotz där Noperschaft gëtt \VUgras{den Hafen} vun de Matrouse ganz vill besicht.

\VUsaut{6.0mm}
\VUfilet{22mm}
